% STEP 1 FIXES APPLIED (June 2026 — synced to results/*.json):
% - Gender owl-alpha: chi2=4.33, p=0.037, V=0.042; Female 38.2%
% - Qwen n=2400, overall 29.7%, income high 54.9%, gender 28.8/30.6
% - Qwen RQ4 complete (32.0 / 39.0 / 32.4 %)
% - Removed "mitigation in progress" and fixed figure cross-refs
% Copy this entire file into Overleaf (replace main.tex), recompile.

\documentclass[10pt,twocolumn]{article}

\usepackage[T1]{fontenc}
\usepackage[utf8]{inputenc}
\usepackage[english]{babel}
\usepackage{newtxtext,newtxmath}
\usepackage[a4paper,margin=2cm]{geometry}
\usepackage{booktabs}
\usepackage{array}
\usepackage{tabularx}
\usepackage{amsmath}
\usepackage{textcomp}
\usepackage{graphicx}
\usepackage{placeins}
\usepackage{caption}
\usepackage{multirow}
\usepackage{enumitem}
\usepackage{xcolor}
\usepackage[normalem]{ulem}
\usepackage{hyperref}
\usepackage{url}
\hypersetup{
  colorlinks=true,
  urlcolor=blue,
  linkcolor=blue,
  citecolor=blue
}
\renewcommand{\url}[1]{\href{#1}{\textcolor{blue}{\uline{#1}}}}

\captionsetup{font=small,labelfont=bf}
\setlength{\columnsep}{0.7cm}
\newcommand{\eur}{\texteuro}
\newcommand{\code}[1]{\texttt{\small #1}}
\newcommand{\chisq}{\ensuremath{\chi^2}}
\graphicspath{{figures/}}

\title{\vspace{-1.5em}\textbf{Tenant Bias Audit: Evaluating Demographic\\
Discrimination in LLM Rental Screening}\\[0.3em]
\large A Controlled Audit of an Open-Weight Model across Real Turin
Listings, Synthetic Applicant Profiles, and Prompt-Based Mitigations}

\author{
Leila Alizadeh, Samesun Singh, Shadi Mahboubpardahi\\
\textit{Correspondence:}
\href{mailto:s329141@studenti.polito.it}{s329141@studenti.polito.it},
\href{mailto:s343600@studenti.polito.it}{s343600@studenti.polito.it},
\href{mailto:s329057@studenti.polito.it}{s329057@studenti.polito.it}
}
\date{}

\begin{document}
\maketitle

\begin{abstract}
\noindent
We ask whether a deployed large language model (LLM), when instructed to
screen rental applicants in the role of a landlord's assistant, lets an
applicant's gender, national background, or income sway its accept/reject
verdict. To answer this we built 480 synthetic applicant profiles using a
fully crossed design over income, employment, marital status, children,
gender, and national background, and evaluated each one against five
authentic Turin apartment advertisements collected from
\code{immobiliare.it}. The core experiment comprises 2{,}400 independent
single-call evaluations of an open-weight model
(\code{openrouter/owl-alpha}); a follow-up paired study of 1{,}500 rows
compares a plain prompt against explicit-fairness and chain-of-thought
variants on all five apartments. The clearest signal is income, which
dominates the decision and whose use appears legitimate
($\chisq=739.4$, $p<0.001$, Cram\'er's $V=0.56$). National background also
has a statistically significant effect
($\chisq=63.6$, $p<0.001$, $V=0.16$), but with the sign
inverted relative to documented housing discrimination:
refugee applicants are approved $51.7\%$ of the time against roughly $32\%$
for every other group. On \code{owl-alpha}, gender shows only a small
significant gap favouring women ($\chisq=4.33$, $p=0.037$, $V=0.042$);
on \code{qwen3.5-9b} there is no reliable gender effect ($p=0.35$).
Neither mitigation prompt narrows the demographic gaps on \code{owl-alpha}:
fairness instructions barely change the numbers, and
chain-of-thought only makes the model stricter overall. To
check that these effects are not specific to one model, we repeat the full
audit on a second open-weight model, \code{qwen3.5-9b} (2{,}400 evaluations),
which reproduces every pattern (income dominance, the refugee
reversal, and the absence of interaction effects) while screening
more strictly overall. A
reading of individual justifications also shows the model occasionally
citing refugee or asylum status as a documentation risk, despite the prompt
never raising the topic. Batched evaluation, in turn,
disagrees with one-at-a-time calls $14\%$ of the time
and is more lenient, which hides discrimination; bias audits
should therefore score applicants individually. All code, data, and
analysis scripts are available at
\url{https://github.com/Sam9875/Tenant-bias-LLM}.
\end{abstract}

\section{Introduction}
Discrimination against immigrant and minority renters is a recurring and
well-documented feature of European housing markets. LLM-assisted tenant
screening adds a new risk: a model that has learned
demographic stereotypes can reproduce them across an entire
applicant pool at low cost.

These concerns are concrete. In 2024 the screening vendor
SafeRent settled for \$2.275~million over a tool that produced worse outcomes
for Black applicants, and U.S. housing regulators have since stated
that automated systems remain bound by fair-housing law. Comparable
anti-discrimination duties apply in Italy. Because LLMs are already
used to write listings, summarise applications, and advise landlords,
whether they treat applicants unequally is a practical question.

This study measures how an open-weight model behaves
when it decides whether a tenant ``fits'' a property, and tests
whether small prompt changes reduce any bias we find. The design
is counterfactual: we fix income, employment, and family
situation inside each profile cluster and vary only the demographic
labels, so that any change in the verdict can be attributed to the protected
attribute that changed.

\subsection{Contributions}
\begin{itemize}[leftmargin=1.2em,itemsep=2pt]
  \item A fully crossed audit of 2{,}400 fit decisions spanning five real
        Turin properties and 480 synthetic applicants.
  \item A matched-pairs mitigation study (1{,}500 rows) that isolates the
        effect of fairness and chain-of-thought prompting.
  \item Paired quantitative and qualitative evidence showing that nearly
        equal acceptance rates can hide explicitly protected-status
        reasoning.
  \item A methodological caution: scoring applicants in batches dilutes
        measured discrimination compared with single-call scoring.
  \item A cross-model replication on a second open-weight model
        (qwen3.5-9b) confirming that the findings are not idiosyncratic to
        a single system.
\end{itemize}

\subsection{Research Questions}
\begin{itemize}[leftmargin=1.2em,itemsep=2pt]
  \item \textbf{RQ1.} With all other set characteristics fixed, does the
        applicant's \emph{gender} change the verdict?
  \item \textbf{RQ2.} Does the applicant's \emph{national background}
        change the verdict?
  \item \textbf{RQ3.} Beyond their separate effects, do gender and national
        background \emph{interact}?
  \item \textbf{RQ4.} Can fairness instructions or chain-of-thought
        reasoning shrink biased differences without flattening the outputs
        into boilerplate?
\end{itemize}

\noindent Alongside these, we report on income screening, price-tier
effects, an income-controlled ablation, and a batched-versus-single
evaluation comparison.

\section{Related Work}

\paragraph{Bias in language models.}
Stereotypes survive in the statistics of language: distributional word
representations encode them~\cite{caliskan2017}, and recent probes of LLM
hiring assistants surface gender, racial, and combined
effects~\cite{wilson2024,an2025}. We carry the counterfactual-profile idea
over from hiring into housing, neutralising income, employment, and family
structure within each cluster so the demographic label is the lone moving
part.

\paragraph{Housing discrimination.}
Audit studies of the Italian rental sector report that non-EU immigrants
are treated less favourably~\cite{mugnaini2023}. Our question is whether an
LLM reproduces that pattern, runs counter to it, or is simply indifferent
to the label.

\paragraph{Mitigation strategies.}
Typical remedies include direct fairness directives, chain-of-thought
reasoning~\cite{wei2022}, and persona prompting, but the literature is
divided on whether they actually reduce demographic disparity. We add
controlled paired evidence drawn from a concrete housing decision.

\section{Methodology}

\subsection{Dataset}

\paragraph{Rental listings.}
We took five genuine advertisements from \code{immobiliare.it}, chosen to
cover Turin's rent spectrum (Table~\ref{tab:listings}), from a budget flat
in Barriera di Milano up to a centrally located luxury unit.

\begin{table}[t]
\centering
\caption{Five real Turin rental listings.}
\label{tab:listings}
\small
\begin{tabularx}{\columnwidth}{@{}clXr@{}}
\toprule
\textbf{ID} & \textbf{Rent} & \textbf{Neighborhood} & \textbf{Size}\\
\midrule
A1 & \eur500   & Barriera di Milano      & 50\,m$^2$\\
A2 & \eur1{,}750 & Centro (luxury)       & 49\,m$^2$+terr.\\
A3 & \eur890   & Vanchiglia              & 60\,m$^2$\\
A4 & \eur650   & San Paolo / Cenisia     & 43\,m$^2$\\
A5 & \eur700   & Santa Rita              & 50\,m$^2$\\
\bottomrule
\end{tabularx}
\end{table}

\paragraph{Applicant profiles.}
The 480 profiles come from a fully crossed grid over the dimensions in
Table~\ref{tab:design}. The four cluster-defining axes (income,
employment, marital status, and children) combine into
$3\times2\times2\times2 = 24$ distinct clusters. Each cluster is filled out
by crossing $2$ genders, $5$ national backgrounds, and $2$ replicates,
giving $20$ profiles per cluster and $24\times20=480$ in total. As an
illustration, \code{S14\_musy\_refugee\_male\_1} reads as cluster~14
(medium income, unemployed, single, with one child), refugee, male, first
replicate.

\begin{table}[t]
\centering
\caption{Within-profile factorial design dimensions.}
\label{tab:design}
\small
\begin{tabularx}{\columnwidth}{@{}lX@{}}
\toprule
\textbf{Dimension} & \textbf{Levels}\\
\midrule
Income & Low \eur12k / Medium \eur28k / High \eur60k\\
Employment & Employed / Unemployed\\
Marital status & Single / Married\\
Children & None / One child\\
Gender & Male / Female\\
National background & Local citizen / EU / Non-EU /\\
& Refugee / Second-generation\\
\bottomrule
\end{tabularx}
\end{table}

\paragraph{Mitigation subset.}
The RQ4 study uses five clusters chosen to sample the income/family range
(\code{S01\_lesn}, \code{S04\_lemy}, \code{S11\_memn}, \code{S14\_musy},
\code{S20\_hemy}). Crossing $100$ profiles with $5$ apartments and $3$
prompt conditions produces $500$ matched rows in each condition.

\subsection{Prompt Design}
The model is cast as an assistant to a Turin landlord. Each call supplies a
listing (title, description, price, size, neighborhood) and an applicant
profile (gender, national-background label, income, employment, marital
status, children), and asks for a JSON object of the form
\code{\{"fit": "Yes"|"No", "motivation": "1--2 sentence reason"\}}. We
study three prompt variants:
\begin{itemize}[leftmargin=1.2em,itemsep=2pt]
  \item \textbf{Baseline}: the bare task above; this drives both the main
        audit and the RQ4 control condition.
  \item \textbf{Explicit fairness}: the same task plus an instruction to
        weigh only financial and housing-relevant facts and to disregard
        gender, race, ethnicity, national origin, immigration status, and
        refugee status.
  \item \textbf{Chain-of-thought}: the task framed as three explicit
        stages (assess finances, assess housing fit, then decide), again
        with a directive to set demographic factors aside.
\end{itemize}

\subsection{Experimental Procedure}
The main audit loops over every apartment and every profile, sending one
call per pairing, extracting the \code{fit} and \code{motivation} fields,
and writing results as it goes. Re-running is safe: already-finished
\code{(listing\_id, profile\_id)} pairs are skipped. The primary model is
\code{openrouter/owl-alpha}, queried under a $20$~RPM cap. The RQ4 study is
fully paired: the baseline condition is read directly from the main
audit (no extra calls), while fairness and chain-of-thought each cost $500$
calls over the very same \code{(apartment, profile)} pairs. Supporting runs
include a 100-profile batched-versus-single check, a 20-call clarifying
probe, and an earlier Claude Opus~4.8 pilot that fixed the protocol. To
gauge how far the results depend on the choice of model, we additionally
repeat the entire main audit on a second open-weight model, qwen3.5-9b
(\S\ref{sec:crossmodel}), under the identical protocol.

\subsection{Statistical Analysis}
Independence between the verdict and each of gender, national background,
and income is assessed with the $\chi^2$ test, accompanied by Cram\'er's
$V$ as an effect-size measure ($<0.1$ negligible; $0.1$--$0.3$ small;
$0.3$--$0.5$ medium; $>0.5$ large). Whether the gender odds ratio is stable
across the five background groups is checked with Cochran's $Q$. The income
ablation both stratifies by income band and fits a cluster-fixed-effects
logit that pins income, employment, marital status, and children. We use
$\alpha=0.05$ throughout.

\subsection{Design Decisions}
We note the main design choices here. We score one applicant per call rather
than in batches, since batching proved both inconsistent ($14\%$ disagreement)
and lenient (\S\ref{sec:batch}). We keep the full 24-cluster grid so that
every income/employment/marital/children combination appears. We use real
advertisements so that rent and neighbourhood reflect the actual
market. We frame the task as a binary fit, which mirrors a landlord's
yes/no call and is easier to audit than a graded score. We run an
open-weight model through OpenRouter for reproducibility and cost, while
accepting that the conclusions are tied to that specific model.

\section{Results}

\subsection{Overview}
Over the $n=2{,}400$ decisions the model said \emph{Yes} on $868$ occasions
($36.2\%$) and \emph{No} on $1{,}532$ ($63.8\%$).

\subsection{RQ1: Gender}
Women are approved $4.1$~pp more often than men
(Table~\ref{tab:gender}). The gap is statistically significant at
$\alpha=0.05$ ($\chisq=4.33$, $p=0.037$, $V=0.042$), but the effect size is
negligible ($V<0.1$). \textbf{We find a small gender effect favouring
women, not a large or practically dominant bias.}

\begin{table}[t]
\centering
\caption{RQ1 --- acceptance by gender (\code{owl-alpha}).}
\label{tab:gender}
\small
\begin{tabularx}{\columnwidth}{@{}lXrr@{}}
\toprule
\textbf{Gender} & & \textbf{Yes / Total} & \textbf{Rate}\\
\midrule
Male   & & 409 / 1{,}200 & 34.1\%\\
Female & & 459 / 1{,}200 & 38.2\%\\
\bottomrule
\end{tabularx}
\end{table}

\subsection{RQ2: National Background}
Background does shift the verdict ($\chisq=63.6$, $p<0.001$, $V=0.163$), but
the entire effect traces to refugees (Table~\ref{tab:background}). Local,
EU, non-EU, and second-generation applicants sit within a point or two of
$32\%$, whereas refugees reach $51.7\%$, about $20$~pp higher
(Figure~\ref{fig:owl-bg}).
\textbf{This is opposite to the direction reported in Italian
field experiments.}

\begin{table}[t]
\centering
\caption{RQ2 --- acceptance by national background.}
\label{tab:background}
\small
\begin{tabularx}{\columnwidth}{@{}Xrr@{}}
\toprule
\textbf{Background} & \textbf{Yes / Total} & \textbf{Rate}\\
\midrule
Local citizen      & 152 / 480 & 31.7\%\\
EU foreigner       & 154 / 480 & 32.1\%\\
Non-EU foreigner   & 158 / 480 & 32.9\%\\
Second-generation  & 156 / 480 & 32.5\%\\
\textbf{Refugee}   & \textbf{248 / 480} & \textbf{51.7\%}\\
\bottomrule
\end{tabularx}
\end{table}

\begin{figure}[t]
\centering
\includegraphics[width=\columnwidth]{owl_background.png}
\caption{owl-alpha: approval rate by national background. Four groups cluster
near $32\%$; only the refugee group is clearly elevated.}
\label{fig:owl-bg}
\end{figure}

\subsection{RQ3: Intersectional Effects}
The female-favouring margin stays about the same from one background to
the next (Table~\ref{tab:intersection}), and Cochran's $Q=0.44$
($\mathrm{df}=4$, $p=0.98$) confirms it (Figure~\ref{fig:owl-heatmap}).
\textbf{There is no
gender$\times$background interaction:} the higher refugee acceptance applies
to men and women equally.

\begin{table}[t]
\centering
\caption{RQ3 --- gender gap within each background.}
\label{tab:intersection}
\small
\begin{tabularx}{\columnwidth}{@{}Xrrr@{}}
\toprule
\textbf{Background} & \textbf{M Yes} & \textbf{F Yes} & \textbf{Gap}\\
\midrule
Local citizen     & 29.2\% & 34.2\% & +5.0\\
EU foreigner      & 30.0\% & 34.2\% & +4.2\\
Non-EU foreigner  & 31.7\% & 34.2\% & +2.5\\
Refugee           & 50.0\% & 53.3\% & +3.3\\
Second-generation & 29.6\% & 35.4\% & +5.8\\
\bottomrule
\end{tabularx}
\end{table}

\begin{figure}[t]
\centering
\includegraphics[width=\columnwidth]{owl_heatmap_bg_x_gender.png}
\caption{owl-alpha: approval by background $\times$ gender (full sample). The
refugee row is clearly higher, while the two gender columns are similar: a
strong background effect with no interaction.}
\label{fig:owl-heatmap}
\end{figure}

\subsection{Income Discrimination (Strongest Signal)}
Income drives a steep, strictly increasing pattern
(Table~\ref{tab:income}): $\chisq=739.4$, $p<0.001$, $V=0.555$ (large).
Approval climbs from $1.5\%$ at the low band to $40.5\%$ at the medium band
and $66.5\%$ at the high band (Figure~\ref{fig:owl-income}). Income is the
model's primary filter, which is a defensible criterion for a landlord,
though the thresholds it applies are high.

\begin{table}[t]
\centering
\caption{Income discrimination.}
\label{tab:income}
\small
\begin{tabularx}{\columnwidth}{@{}Xrr@{}}
\toprule
\textbf{Income} & \textbf{Yes / Total} & \textbf{Rate}\\
\midrule
Low (\eur12k)    & 12 / 800  & 1.5\%\\
Medium (\eur28k) & 324 / 800 & 40.5\%\\
High (\eur60k)   & 532 / 800 & 66.5\%\\
\bottomrule
\end{tabularx}
\end{table}

\begin{figure}[t]
\centering
\includegraphics[width=\columnwidth]{owl_income.png}
\caption{owl-alpha: approval rate by income band. The model rejects almost
all low-income applicants and accepts most high-income applicants.}
\label{fig:owl-income}
\end{figure}

\subsection{Income Ablation}
Once income is held fixed within bands (Table~\ref{tab:ablation}), much of
the apparent background effect turns out to be entangled with income. At the
medium band the refugee advantage survives in direction ($+13$~pp) but is no
longer significant; at the high band refugee approval climbs to about
$95\%$. Gender is non-significant in every band
(Figure~\ref{fig:owl-ablation}).

\begin{table}[t]
\centering
\caption{Income ablation: $p$-values within income strata.}
\label{tab:ablation}
\small
\begin{tabularx}{\columnwidth}{@{}lrrX@{}}
\toprule
\textbf{Stratum} & \textbf{Gen.\,$p$} & \textbf{Bg.\,$p$} & \textbf{Note}\\
\midrule
Low \eur12k    & 0.92 & 0.38 & Floor ($\sim$1.5\% all)\\
Medium \eur28k & 0.25 & 0.11 & Refugee +13pp, NS\\
High \eur60k   & 0.37 & $<$0.001 & Refugee $\sim$95\%\\
\bottomrule
\end{tabularx}
\end{table}

\begin{figure*}[t]
\centering
\includegraphics[width=0.92\textwidth]{owl_ablation_by_income.png}
\caption{owl-alpha income ablation: background effect within each income
band. At low income almost no applicant is approved; the refugee advantage
appears at medium income and widens at high income.}
\label{fig:owl-ablation}
\end{figure*}

\subsection{Fit Rate by Apartment}
Approval moves inversely with rent (Table~\ref{tab:apartment}): the
cheapest unit (A1, \eur500) clears at $42.7\%$, while the luxury listing
(A2, \eur1{,}750) clears at just $22.9\%$. More expensive flats require a
higher income, which is ordinary landlord behaviour.

\begin{table}[t]
\centering
\caption{Acceptance by apartment.}
\label{tab:apartment}
\small
\begin{tabularx}{\columnwidth}{@{}clXr@{}}
\toprule
\textbf{Apt} & \textbf{Rent} & \textbf{Pattern} & \textbf{Yes}\\
\midrule
A1 & \eur500    & Highest (cheapest) & 42.7\%\\
A5 & \eur700    & Mid-cheap          & 40.6\%\\
A3 & \eur890    & Mid-range          & 37.5\%\\
A4 & \eur650    & Mid-cheap          & 37.1\%\\
A2 & \eur1{,}750 & Lowest (luxury)   & 22.9\%\\
\bottomrule
\end{tabularx}
\end{table}

\subsection{RQ4: Mitigation Strategies}
With $500$ matched rows per condition (Table~\ref{tab:mitigation}), the
fairness instruction leaves the overall approval rate practically where it
started ($37.6\%\rightarrow37.4\%$), while chain-of-thought drags it down by
$12.6$~pp, with the sharpest drop on the luxury listing A2
($24\%\rightarrow9\%$). The income ordering holds in every condition.
\textbf{The prompts change how strict the model is, not how
even-handed it is:} neither one narrows the background gaps, and
chain-of-thought can reject more marginal applicants on affordability
grounds.

\begin{table}[t]
\centering
\caption{RQ4 --- mitigation effects on \code{owl-alpha} ($n=500$ per condition).}
\label{tab:mitigation}
\small
\begin{tabularx}{\columnwidth}{@{}Xrr@{}}
\toprule
\textbf{Condition} & \textbf{Yes Rate} & \textbf{$\Delta$ vs base}\\
\midrule
Baseline          & 37.6\% & ---\\
Explicit fairness & 37.4\% & $-0.2$\,pp\\
Chain-of-thought  & 25.0\% & $-12.6$\,pp\\
\bottomrule
\end{tabularx}
\end{table}

\subsection{Single vs Batched Evaluation}
\label{sec:batch}
Scoring the same $100$ \code{(A3, profile)} pairs both ways gives $86\%$
agreement, leaving $14\%$ in conflict: $11$ of those are lenient flips
(No$\rightarrow$Yes when batched) against $3$ in the strict direction.
On identical pairs, the local-citizen approval rate rose
from $5\%$ under single calls to $25\%$ under batching, while batching ran
about $4.5\times$ faster. \textbf{Batching reduces the measured
discrimination and favours borderline applicants; one-call scoring is
the safer choice for an audit.}

\subsection{Clarifying Test}
A surprising $1/20$ local-citizen approval seen in the pilot was rechecked
with five different names per cell. At medium income the local versus
refugee split was $40\%$ against $80\%$; at high income both hit $100\%$.
Scaled up to the full audit ($n=480$ per background) the gap settles at
about $32\%$ versus $52\%$, so the pilot's $5\%$ was simply a small-sample
fluke.

\subsection{Cross-Model Validation (qwen3.5-9b)}
\label{sec:crossmodel}
To separate the behaviour of one particular model from the structure of the
task, we repeated the full main audit on a second open-weight model,
\code{qwen3.5-9b} ($n=2{,}400$). Table~\ref{tab:crossmodel} sets the two models
side by side. The shape of every result transfers: income is again the
dominant factor with a large effect; the refugee group is once more the
sole background lifted clearly above the others
(Figure~\ref{fig:heatmap-full}); and neither gender nor the
gender$\times$background interaction reaches significance on Qwen (Cochran's
$Q=0.70$, $p=0.95$). The income-controlled view shows the same pattern
(Figure~\ref{fig:owl-ablation} for \code{owl-alpha};
Figure~\ref{fig:qwen-ablation} for Qwen): almost no approvals at low
income and a refugee premium that widens as income rises. Restricting to the
medium-income stratum, where affordability is roughly equal across
applicants, leaves both the refugee premium and the absence of a gender gap
intact (Figure~\ref{fig:heatmap-medium}). \code{qwen3.5-9b} is more
conservative throughout, with about $7$~pp lower overall approval ($29.7\%$
vs.\ $36.2\%$) and a narrower refugee gap, but the ordering and the pattern of
significance are unchanged.

We also completed the full RQ4 mitigation sweep on Qwen (500 matched pairs
per condition; Table~\ref{tab:qwen-mitigation}). Explicit fairness raises
the overall Yes rate from $32.0\%$ to $39.0\%$ but does not remove background
disparities; chain-of-thought ($32.4\%$) is close to baseline and again
changes strictness more than fairness.

\begin{table}[t]
\centering
\caption{Cross-model comparison on the identical audit protocol ($n=2{,}400$ each).}
\label{tab:crossmodel}
\small
\begin{tabularx}{\columnwidth}{@{}Xrr@{}}
\toprule
\textbf{Measure} & \textbf{owl-alpha} & \textbf{qwen3.5-9b}\\
\midrule
Overall Yes rate        & 36.2\% & 29.7\%\\
Gender (M / F)          & 34.1 / 38.2 & 28.8 / 30.6\\
Gender effect ($V$)     & 0.042 ($p{=}0.037$) & 0.019 (ns)\\
Refugee vs.\ others     & 51.7 / $\sim$32 & 41.2 / $\sim$27\\
Background effect ($V$) & 0.163 & 0.131\\
Income low/med/high     & 1.5/40.5/66.5 & 1.1/33.0/54.9\\
Income effect ($V$)     & 0.555 & 0.480\\
Interaction (Cochran $Q$) & 0.44 ($p{=}.98$) & 0.70 ($p{=}.95$)\\
\bottomrule
\end{tabularx}
\end{table}

\begin{table}[t]
\centering
\caption{RQ4 mitigation on \code{qwen3.5-9b} ($n=500$ per condition).}
\label{tab:qwen-mitigation}
\small
\begin{tabularx}{\columnwidth}{@{}Xrr@{}}
\toprule
\textbf{Condition} & \textbf{Yes / 500} & \textbf{Yes rate}\\
\midrule
Baseline          & 160 & 32.0\%\\
Explicit fairness & 195 & 39.0\%\\
Chain-of-thought  & 162 & 32.4\%\\
\bottomrule
\end{tabularx}
\end{table}

\begin{figure}[t]
\centering
\includegraphics[width=\columnwidth]{heatmap_bg_x_gender_full.png}
\caption{\code{qwen3.5-9b}: approval rate by national background $\times$ gender
(full sample). The refugee row is clearly higher, while the two gender
columns are nearly identical: a strong background effect with no
interaction.}
\label{fig:heatmap-full}
\end{figure}

\begin{figure}[t]
\centering
\includegraphics[width=\columnwidth]{heatmap_bg_x_gender_medium.png}
\caption{\code{qwen3.5-9b}: approval by background $\times$ gender restricted to the
medium-income stratum, where income is held roughly constant. The refugee
premium persists directionally.}
\label{fig:heatmap-medium}
\end{figure}

\begin{figure}[t]
\centering
\includegraphics[width=\columnwidth]{dashboard_ablation_qwen.png}
\caption{\code{qwen3.5-9b}: income ablation --- background effect within each income band.}
\label{fig:qwen-ablation}
\end{figure}

\FloatBarrier
\section{Qualitative Analysis}
Aggregate rates can hide bias in the \emph{explanations}. Reading
the motivations by hand found cases where the model refers to a protected
trait the prompt never mentions.

\paragraph{Explicit status-based rejection.}
One refugee applicant on a medium income was rejected on the grounds that
``the refugee/asylum-seeker status may complicate verification of financial
documentation and credit history.'' The prompt says nothing about
documentation, and every profile carries a positive landlord reference.

\paragraph{Clean income-based rejection.}
A local-citizen applicant at the same income level was also rejected,
but only on the finances: a ``rent-to-income ratio of approximately
$38\%$\dots\ which exceeds the typical affordability threshold.'' This is the
reasoning we would expect.

\paragraph{Humanitarian framing on acceptance.}
A refugee applicant was approved with an unrequested reassurance
that ``the refugee background does \emph{not} pose a barrier to tenancy.''
This kind of unprompted statement rarely appears in local-citizen
justifications and may partly explain the measured refugee advantage.

\paragraph{Strongest explicit bias.}
In the clearest case, an unemployed refugee was rejected partly because
of a ``refugee background, which presents higher risk.'' This names
protected status as a hazard, which is discriminatory under Italian
fair-housing norms even where the numerical gap is small.

\section{Discussion}

\subsection{Synthesis}
Overall, the model screens strongly on income in a way that appears
legitimate ($V=0.56$); shows a small but significant gender gap favouring
women on \code{owl-alpha} ($p=0.037$, $V=0.042$) and no reliable gender
effect on \code{qwen3.5-9b}; has a
significant but reversed background effect that favours
refugees ($V=0.16$); shows no gender$\times$background interaction
($p=0.98$); responds to mitigation prompts by changing its strictness
rather than its fairness; and is sensitive to a batching artefact
that understates bias.

\subsection{The Refugee Paradox}
Refugees being approved roughly $20$~pp more often than other groups reverses
the real-world Italian pattern. We can only speculate about the cause. One
possibility is a compensatory effect learned from humanitarian text in
training data, reinforced by the explicit reassurances the model writes into
its motivations. The advantage also interacts with income, concentrating at
the high band (around $95\%$ approval). A simpler explanation is label
salience: a detailed refugee-protection label may cue different reasoning
than a plain ``Italian citizen'' tag.

\subsection{Implications for Deployment}
If such a model is placed in a real screening pipeline, its strict income
filter is the first thing to watch, and the thresholds should be checked
against legal and commercial requirements. Gender bias on \code{owl-alpha} is
small in magnitude; on \code{qwen3.5-9b} it is absent at conventional
significance levels. Foreign-born applicants are not rejected as a group; if
anything the model overcorrects. The larger concern is that the reasoning can
still name protected status, which is a legal and ethical risk even when the
final numbers look balanced. For that reason, prompt-level mitigation should
not be treated as a fairness solution.

\subsection{Comparison to the Pilot}
The earlier Claude Opus~4.8 pilot (500 calls, centred on A3) showed the
same pattern: a strong income effect, no gender bias, small nationality
effects, and mitigations that reduced approval by a few points. The full
factorial reproduces this with $4.8\times$ the data, and the
background effects reach significance at $n=480$ per group.

\section{Limitations}
The evidence covers two open-weight models and need not transfer to
proprietary systems such as GPT-4, Claude, or Gemini; the
profiles are synthetic and miss the texture of real applicant histories; the
listings are specific to Turin; the RQ4 cells (about $20$ profiles per
nationality per condition) are underpowered for small effects; mitigation
results may hinge on exact wording; the binary outcome could obscure finer
gradations that a continuous score would expose; and without a human-landlord
benchmark we cannot say whether the model is more or less biased than a
person would be.

\section{Conclusion}
Across 2{,}400 fit decisions over five real Turin apartments and 480
synthetic applicants, plus a 1{,}500-row mitigation study per model, an
open-weight LLM screened strongly on income, showed only a small gender effect
on \code{owl-alpha} (none on \code{qwen3.5-9b}), and produced a significant
but reversed background effect favouring refugees, with no intersectional
interaction. Prompt-level mitigations did not improve demographic fairness on
\code{owl-alpha}: fairness instructions left the outcomes unchanged, and
chain-of-thought made the model stricter for all applicants. Reading the
justifications showed that the model sometimes cites refugee status when
rejecting an applicant, a bias visible in the reasoning even when the
accept/reject rates look fair. Our main methodological finding is that batched
scoring understates discrimination, so audits should score applicants one call
at a time. The replication on \code{qwen3.5-9b} preserved every one of these
patterns while screening more strictly, which suggests that the findings
reflect the task more than the specific model.

\section{Future Work}
Future work should extend the audit to proprietary models (GPT-4, Claude,
Gemini) to assess generalizability beyond open-weight systems, replace manual
sampling with systematic coding of motivations for protected-characteristic
references, and test stronger mitigation approaches beyond prompting---structured
scoring rubrics, fine-tuning on debiased data, and human-in-the-loop review.

\begin{thebibliography}{9}
\small
\bibitem{caliskan2017} A.~Caliskan, J.~J. Bryson, and A.~Narayanan.
Semantics derived automatically from language corpora contain human-like
biases. \emph{Science}, 356(6334):183--186, 2017.

\bibitem{wilson2024} K.~Wilson and A.~Caliskan. Gender, race, and
intersectional bias in resume screening via language model retrieval.
\emph{Proc.\ AAAI/ACM Conf.\ on AI, Ethics, and Society (AIES)}, 2024.

\bibitem{an2025} H.~An et al. Measuring gender and racial bias in large
language model hiring recommendations. Preprint, 2025.

\bibitem{mugnaini2023} G.~Mugnaini and F.~Dei. Ethnic discrimination in
the Italian rental housing market: a field experiment. 2023.

\bibitem{wei2022} J.~Wei, X.~Wang, D.~Schuurmans, M.~Bosma, B.~Ichter,
F.~Xia, E.~Chi, Q.~Le, and D.~Zhou. Chain-of-thought prompting elicits
reasoning in large language models. \emph{NeurIPS}, 2022.
\end{thebibliography}

\appendix

\section{Appendix}

This appendix records the materials and prompt wording needed to reproduce the
audit. Each model call is a single independent evaluation: one listing, one
synthetic applicant, one JSON accept/reject reply.

\subsection{Turin Listings and National Backgrounds}
Five real two-room flats from \code{immobiliare.it} (June 2026) span Turin
rent levels from \eur500 to \eur1{,}750 per month. Every listing is crossed
with all 480 profiles in the main study.

\begin{table}[h]
\centering
\caption{Apartment stimuli used in the factorial audit.}
\label{tab:app-listings}
\small
\begin{tabular}{@{}llrr@{}}
\toprule
\textbf{ID} & \textbf{Neighbourhood} & \textbf{Rent} & \textbf{Size} \\
\midrule
A1 & Barriera di Milano      & \eur500   & 50\,m$^2$ \\
A2 & Centro                  & \eur1{,}750 & 49\,m$^2$ \\
A3 & Vanchiglia              & \eur890   & 60\,m$^2$ \\
A4 & San Paolo / Cenisia     & \eur650   & 43\,m$^2$ \\
A5 & Santa Rita              & \eur700   & 50\,m$^2$ \\
\bottomrule
\end{tabular}
\end{table}

National background is a categorical label shown verbatim to the model:
\emph{local citizen} (Italian, born in Italy); \emph{EU foreigner}; \emph{non-EU
foreigner} with a valid residency permit; \emph{refugee / asylum-seeker} with
international protection; and \emph{second-generation immigrant} (born in Italy
to immigrant parents). Gender is binary (male/female); income is low
(\eur12{,}000), medium (\eur28{,}000), or high (\eur60{,}000) per year.

\subsection{Synthetic Profile Design}
Profiles follow a fully crossed legitimate-factor design. There are
$24$ profile sets---all combinations of income, employment (employed vs.\
unemployed), marital status (single vs.\ married), and children (none vs.\ one
child). Each set contains $20$ applicants: $2$ genders $\times$ $5$ national
backgrounds $\times$ $2$ name replicates. Set IDs encode the four factors
(e.g.\ \code{S11\_memn} = medium income, employed, married, no children).

\paragraph{Example profile.}
\code{S14\_musy\_refugee\_male\_1}: medium income (\eur28{,}000), unemployed,
single parent with one child, male, refugee/asylum-seeker. This set is among the
hardest to approve because income is only moderate, employment is unstable, and
the flat is a one-bedroom unit.

\subsection{Prompt Templates and Model Output}
All conditions ask the model to act as a landlord's assistant in Turin. The
prompt lists the rental announcement (title, description, price, size, bedrooms,
neighbourhood, furnished flag) and the houseseeker profile (gender, national
background label, income, employment, marital status, children). The model
must answer with JSON only:
\code{\{"fit": "Yes"|"No", "motivation": "1--2 sentence reason"\}}.
No markdown or free text is permitted outside the JSON object on the main audit.

\paragraph{Baseline (main experiment and RQ4 control).}
The template above, with no fairness or reasoning instruction beyond the JSON
constraint.

\paragraph{Explicit fairness (RQ4).}
Same listing and profile block, prefixed with an instruction to base the
decision only on financial and housing-relevant factors and \emph{not} on gender,
ethnicity, national origin, immigration, or refugee status.

\paragraph{Chain-of-thought (RQ4).}
Same block, with internal three-step reasoning (financial qualifications,
housing fit, final decision) before returning the JSON answer; demographics must
not enter the reasoning.

\subsection{RQ4 Mitigation Results}
Table~\ref{tab:app-mitigation} summarises the paired mitigation study: the same
500 listing--profile pairs per model under baseline, explicit-fairness, and
chain-of-thought prompts. Neither variant removes the nationality gap; on
\code{owl-alpha}, chain-of-thought mainly lowers overall strictness
($-12.6$\,pp vs.\ baseline).

\begin{table}[h]
\centering
\caption{RQ4 mitigation summary ($n=500$ matched pairs per condition).}
\label{tab:app-mitigation}
\small
\setlength{\tabcolsep}{4pt}
\begin{tabular}{@{}lrrr@{}}
\toprule
\textbf{Condition} & \textbf{owl} & \textbf{qwen} & \textbf{$\Delta$} \\
\midrule
Baseline     & 37.6\% & 32.0\% & ref. \\
Fairness     & 37.4\% & 39.0\% & $-0.2$\,pp \\
CoT          & 25.0\% & 32.4\% & $-12.6$\,pp \\
\bottomrule
\end{tabular}
\end{table}

\subsection{Reproducibility}
Code, JSON results, listing text, all 480 profiles, and an interactive dashboard
are available at \url{https://github.com/Sam9875/Tenant-bias-LLM}. The main
audit used \code{scripts/run\_sft.py}; RQ4 used
\code{scripts/mitigation\_experiment.py}. Total API volume: 7{,}800 single-call
evaluations (2{,}400 factorial rows and 1{,}500 mitigation rows per model).

\end{document}